% BEGIN VOL07-EXPANSION VII15
\section{Supplementary worked examples}
\label{sec:vii15-pedagogy}

\begin{example}[Metric of a graph]\label{ex:vii15-ped-01}
For \(X(u,v)=(u,v,f(u,v))\), the first fundamental form has
\[
E=1+f_u^2,\qquad F=f_uf_v,\qquad G=1+f_v^2.
\]
It converts parameter velocities into intrinsic lengths and angles on the surface.
\end{example}

\begin{example}[Plane and cylinder share the same local metric]\label{ex:vii15-ped-02}
The plane patch \(X(u,v)=(u,v,0)\) and the cylinder patch \(Y(u,v)=(\cos u,\sin u,v)\) both have
\[
I=du^2+dv^2.
\]
Yet their second fundamental forms differ.  This is the basic example separating intrinsic metric geometry from extrinsic bending.
\end{example}

\begin{example}[Second fundamental form of a sphere]\label{ex:vii15-ped-03}
With outward unit normal \(N=X/R\) on a sphere of radius \(R\), the second fundamental form is proportional to the first.  Depending on the sign convention for the shape operator, one writes \(II=\pm R^{-1}I\).  The proportionality, not the convention-dependent sign, expresses that every direction is principal.
\end{example}

\section{Graded supplementary exercises}
The following exercises add a deliberate progression from concrete calculation to structural proof, hypothesis testing, application, and synthesis.

\subsection*{Standard computations and constructions}

\begin{exercise}[Plane first form]\label{exr:vii15-ped-01}
For \(X(u,v)=(u,v,0)\), compute \(E,F,G\).
\end{exercise}
\begin{hint}
Take dot products of \(X_u,X_v\).
\end{hint}
\begin{solution}
\(E=1,\ F=0,\ G=1\).
\end{solution}

\begin{exercise}[Cylinder first form]\label{exr:vii15-ped-02}
For \(X(u,v)=(R\cos u,R\sin u,v)\), compute \(E,F,G\).
\end{exercise}
\begin{hint}
Differentiate and take dot products.
\end{hint}
\begin{solution}
\(E=R^2,\ F=0,\ G=1\).
\end{solution}

\begin{exercise}[Graph first form]\label{exr:vii15-ped-03}
For \(f(u,v)=u^2+v^2\), compute \(E,F,G\).
\end{exercise}
\begin{hint}
Use \(f_u=2u,\ f_v=2v\).
\end{hint}
\begin{solution}
\(E=1+4u^2,\ F=4uv,\ G=1+4v^2\).
\end{solution}

\begin{exercise}[Area density]\label{exr:vii15-ped-04}
Express \(\|X_u\times X_v\|\) in terms of \(E,F,G\).
\end{exercise}
\begin{hint}
Use the Gram determinant.
\end{hint}
\begin{solution}
\(\|X_u\times X_v\|=\sqrt{EG-F^2}\).
\end{solution}

\begin{exercise}[Angle formula]\label{exr:vii15-ped-05}
How is the angle between tangent vectors with coordinate columns \(a,b\) computed from the first fundamental form matrix \(g\)?
\end{exercise}
\begin{hint}
Use the induced inner product.
\end{hint}
\begin{solution}
\(\cos\theta=(a^Tgb)/\sqrt{(a^Tga)(b^Tgb)}\).
\end{solution}

\subsection*{Proofs}

\begin{exercise}[Positive definiteness]\label{exr:vii15-ped-06}
Prove \(EG-F^2>0\) for a regular surface patch.
\end{exercise}
\begin{hint}
Use the Gram determinant of two independent tangent vectors.
\end{hint}
\begin{solution}
\(EG-F^2=\|X_u\times X_v\|^2>0\) because regularity makes the cross product nonzero.
\end{solution}

\begin{exercise}[First form invariance]\label{exr:vii15-ped-07}
Explain why the first fundamental form is independent of parametrization.
\end{exercise}
\begin{hint}
It is the ambient Euclidean inner product restricted to the tangent plane.
\end{hint}
\begin{solution}
Changing coordinates only changes the matrix representing the same bilinear form; the geometric inner product itself is unchanged.
\end{solution}

\begin{exercise}[Normal reversal]\label{exr:vii15-ped-08}
Prove that reversing the chosen unit normal reverses the second fundamental form but leaves the first unchanged.
\end{exercise}
\begin{hint}
Track where the normal enters the definitions.
\end{hint}
\begin{solution}
The first form uses only tangent dot products.  The second form is linear in the normal (or shape operator), so \(N\mapsto-N\) multiplies it by \(-1\).
\end{solution}

\begin{exercise}[Isometry preserves first form]\label{exr:vii15-ped-09}
Prove a local surface isometry preserves lengths of tangent vectors.
\end{exercise}
\begin{hint}
Use equality of pulled-back first fundamental forms.
\end{hint}
\begin{solution}
If \(\Phi^*I_2=I_1\), then \(I_2(d\Phi v,d\Phi v)=I_1(v,v)\) for every tangent vector \(v\).
\end{solution}

\subsection*{Counterexamples and hypothesis tests}

\begin{exercise}[Same first form means same embedding]\label{exr:vii15-ped-10}
Test with a plane and cylinder.
\end{exercise}
\begin{hint}
Compare \(I\) and extrinsic shape.
\end{hint}
\begin{solution}
False.  They can have the same local metric while one is flat as an embedded plane and the other bends in \(\mathbb R^3\).
\end{solution}

\begin{exercise}[Second form is orientation-independent]\label{exr:vii15-ped-11}
Test under \(N\mapsto-N\).
\end{exercise}
\begin{hint}
The second form uses the normal.
\end{hint}
\begin{solution}
False.  Its sign reverses with the unit normal.
\end{solution}

\begin{exercise}[\(F=0\) means unit coordinates]\label{exr:vii15-ped-12}
Test the meaning of \(F=0\).
\end{exercise}
\begin{hint}
Check \(E\) and \(G\) separately.
\end{hint}
\begin{solution}
False.  \(F=0\) means coordinate directions are orthogonal, but their lengths are \(\sqrt E\) and \(\sqrt G\), not necessarily \(1\).
\end{solution}

\subsection*{Applications and investigations}

\begin{exercise}[Surface area]\label{exr:vii15-ped-13}
Explain how the first fundamental form determines area.
\end{exercise}
\begin{hint}
Use the Gram determinant.
\end{hint}
\begin{solution}
The area element in parameters is \(\sqrt{EG-F^2}\,du\,dv\).
\end{solution}

\begin{exercise}[Intrinsic path length]\label{exr:vii15-ped-14}
How does \(I\) compute the length of a surface curve \(X(u(t),v(t))\)?
\end{exercise}
\begin{hint}
Insert the parameter velocity into the quadratic form.
\end{hint}
\begin{solution}
The speed squared is \(E\dot u^2+2F\dot u\dot v+G\dot v^2\), whose square root is integrated along the curve.
\end{solution}

\subsection*{Challenge problems}

\begin{exercise}[Plane-cylinder local isometry]\label{exr:vii15-ped-15}
Construct an explicit local isometry from the plane to a unit cylinder.
\end{exercise}
\begin{hint}
Use \((u,v)\mapsto(\cos u,\sin u,v)\).
\end{hint}
\begin{solution}
The differential sends \(\partial_u,\partial_v\) to orthonormal vectors, so the pulled-back cylinder metric is \(du^2+dv^2\).
\end{solution}

\begin{exercise}[Umbilic criterion via forms]\label{exr:vii15-ped-16}
Explain what it means for \(II\) to be proportional to \(I\) at a point.
\end{exercise}
\begin{hint}
Relate the shape operator to the identity.
\end{hint}
\begin{solution}
It means every tangent direction has the same normal curvature; equivalently the shape operator is a scalar multiple of the identity, so the point is umbilic.
\end{solution}

% END VOL07-EXPANSION VII15
